\documentclass[12pt,letterpaper]{article}
\usepackage[utf8]{inputenc}
\usepackage[T1]{fontenc}
\usepackage{times}
\usepackage[margin=0.5in,top=0.4in,bottom=0.4in]{geometry}
\usepackage{hyperref}
\usepackage{xcolor}
\usepackage{enumitem}
\usepackage{titlesec}

\hypersetup{colorlinks=true,linkcolor=blue!60!black,urlcolor=blue!60!black,citecolor=blue!60!black}

\setlength{\parindent}{0pt}
\setlength{\parskip}{2pt}
\setlist{nosep,leftmargin=*}

\titleformat{\section}{\normalsize\bfseries\scshape}{}{0pt}{}[\vspace{-2pt}\rule{\linewidth}{0.4pt}]
\titleformat{name=\section,numberless}{\large\bfseries\color{blue!60!black}}{}{0pt}{}[\vspace{-2pt}\rule{\linewidth}{0.6pt}]
\titlespacing{\section}{0pt}{8pt}{4pt}

\pagestyle{empty}

\begin{document}

{\footnotesize\textbf{Arxiv Digest --- 2026-08-23} \hfill 7 papers}

\smallskip

\section*{Multilingual NLP}

{\small\textbf{Shared Circuits for Shared Grammar: Tracing Subject-Verb Agreement Across Languages}} {\scriptsize[\href{https://arxiv.org/abs/2608.18545}{arxiv}]}\\
{\scriptsize Isabella Gidi, Antonio Almudévar, Core Francisco Park, Naomi Saphra, Ricard Marxer}\\

{\small\textit{Multilingual LLMs often reuse the same subject--verb agreement circuit across languages, especially when the language has overt agreement morphology.}}\\[1pt]
{\footnotesize Mechanistically tests whether cross-lingual grammatical generalization uses shared vs language-specific machinery by tracing a causal agreement circuit in 29 languages and quantifying head overlap as a function of overt person/number marking. Use activation patching to identify attention heads that causally restore correct present-tense verb inflection in corrupted sentences; compare implicated head sets and their attention targets across languages to separate shared head locations from shared function. Agreement-head overlap is much higher in languages with overt verb inflection; sharing is strongest when evaluation isolates the inflectional contrast itself. English behaves like inflected languages specifically in contexts with overt agreement (-s), and many shared heads show similar attention patterns, indicating genuinely shared computation.}

\medskip

{\small\textbf{Auditing Cross-Lingual Fairness in Language Model Watermarking}} {\scriptsize[\href{https://arxiv.org/abs/2608.20047}{arxiv}]}\\
{\scriptsize Alexander Nemecek, Osama Zafar, Debargha Ganguly, Vikash Singh, Vipin Chaudhary, Erman Ayday}\\

{\small\textit{Watermarking that looks reliable in English can be unfair or brittle cross-lingually, largely due to thresholding and metric choices that don't transfer across languages.}}\\[1pt]
{\footnotesize Introduces an audit framework that disentangles watermark method failures from evaluation-design failures, and attributes cross-lingual disparities to structural (typology/script/tokenization) factors via family-level decomposition. (1) Calibrate detection thresholds per language/context; (2) add a threshold-independent separability measure; (3) assess quality with three paradigms (distributional, paired semantic, reference-perplexity); (4) quantify and decompose disparity using generalized-entropy over typological families. Across 6 schemes, 3 generators, 11 languages, and base vs instruction-tuned settings, detection conclusions often flip under proper calibration; the separability metric reveals when detectors truly fail vs are mis-thresholded. Quality impacts vary by metric and language, and most disparity is between typological families, implying systematic cross-lingual unfairness.}

\medskip

{\small\textbf{Institutional Books - Enriched Text: A customizable multilingual open-source pipeline for denoising, deduplicating, and annotating OCR text at scale}} {\scriptsize[\href{https://arxiv.org/abs/2608.19026}{arxiv}]}\\
{\scriptsize David Lowry-Duda, Matteo Cargnelutti, Catherine Brobston, Salwa Ismail, Greg Leppert, Amanda Watson, Jonathan Zittrain}\\

{\small\textit{Instead of one irreversible ``cleaned OCR'' dataset, this work ships normalized OCR text plus rich per-paragraph annotations so users can apply their own filtering, language, and dedup rules.}}\\[1pt]
{\footnotesize Proposes ``Enriched Text'' as a non-destructive packaging of large OCR corpora---separating measurement from deletion---and releases IB-HL-ET plus an open pipeline that scales to \textasciitilde{}250 languages while preserving provenance and quality signals. Normalize/segment OCR books and attach HTML-like paragraph annotations: language ID (incl. multilingual), endmatter labels, dedup cluster IDs, and bits-per-byte compressibility as a noise proxy---stored as tags/attributes rather than used for hard filtering. Demonstrates feasibility at library scale: 983,003 volumes, \textasciitilde{}1.39B annotated paragraphs, \textasciitilde{}217B tokens. The artifact enables reproducible slicing by language/quality/duplication without rerunning processing or accepting one-size-fits-all filtering that can destroy valuable content.}

\medskip

{\small\textbf{HealMed: Multilingual Evaluation of Large Language Models in Medicine}} {\scriptsize[\href{https://arxiv.org/abs/2608.19981}{arxiv}]}\\
{\scriptsize Yingjian Chen, Fan Gao, Sherry T. Tong, Haoyu Zhang, Aosong Feng, Kevin W. Jin, Xing Wu, Jinghui Lu, Abdul Samad, Akbar Faruqi, Cesar Caraballo, Cibele Brandão,  Dhruva,  Gupta, Eunji Jeon, Gabriel Madera-Santiago, Geon Lee, Hugo Toshio Itikawa, Insook Cho, Isabelli Martins, Isarar Siddique, Israr Ahmed, Jihyo Kwak, Kanyakorn Veerakanjana, Luis Guilherme Cardoso, Minjin Kim, Piyalitt Ittichaiwong, Renee Dua, Santiago Gudiño-Rosales, Xiujie Chen, Zeo Lapalus, Zixin Xu, Michihiro Yasunaga, Rex Ying, Heuiseok Lim, Jaewoo Kang, Chanjun Park, Hang Jiang, Ethan Goh, Hyunjae Kim, Edison Marrese-Taylor, Yusuke Iwasawa, Yutaka Matsuo, Qingyu Chen, Irene Li}\\

{\small\textit{Multilingual medical LLM evaluation is highly confounded by translation quality; HealMed uses expert-revised translations to make cross-language results credible.}}\\[1pt]
{\footnotesize Builds HealMed, a 9-language medical benchmark where every translated item is independently reviewed/revised by bilingual medical experts, showing that non-expert translation can systematically distort measured ``multilingual medical ability.'' Create 1,000 examples per language from 9 medical datasets mapped into MCQA, NLI, and open QA. Translate from English, then run a two-expert audit/revision to fix meaning, terminology, pragmatics, and artifacts (e.g., negation shifts, unnatural cues). Evaluate diverse proprietary/open/medical LLMs and compare to unrevised translations. Performance drops are real and often largest in lower-resource languages, but vary strongly by model and task. Expert revision can raise or lower scores, proving translation artifacts can both penalize and inflate results; thus benchmarks without domain-expert translation control can mis-rank models for non-English clinical use.}

\medskip

{\small\textbf{TranslatePsy-AfriSLM: High-Quality Data Scaling For Low-Resource Machine Translation}} {\scriptsize[\href{https://arxiv.org/abs/2608.18655}{arxiv}]}\\
{\scriptsize Milan Gritta, Patrik Lambert, Jihye Back, Amril Nazir}\\

{\small\textit{For low-resource African MT, aggressive quality filtering plus targeted synthetic data lets sub‑1B models beat much larger LLMs.}}\\[1pt]
{\footnotesize Releases TranslatePsy-AfriSLM for 19 Sub-Saharan African languages and shows a practical recipe: unify quality estimation across sources, discard most noisy pairs, and rely heavily on filtered synthetic parallel data to maximize quality per compute. Apply a single QE-based filter to both mined/curated parallel and synthetic data, removing up to 96\% of tokens; train/fine-tune small MT models on the remaining high-signal mixture and compare against much larger general-purpose models. Extreme filtering is beneficial (dropping up to 96\% tokens does not hurt quality), and filtered synthetic data gives the best efficiency. Fine-tuned models as small as 0.8B outperform much larger systems including TranslateGemma-27B and Qwen3.5-122B-A10B.}

\medskip


\section*{Multilingual Speech Processing}

{\small\textbf{Understanding Multilingual Medical ASR Adaptation Through Layer-Wise Analysis}} {\scriptsize[\href{https://arxiv.org/abs/2608.18825}{arxiv}]}\\
{\scriptsize Souranil Kahali, Rituparna Bose, Abner Hernandez, Tomas Arias-Vergara, Andreas Maier, Ning Ma, Paula Andrea Perez-Toro}\\

{\small\textit{In multilingual medical Whisper adaptation, domain tuning drives the main internal representation shift; later multilingual tuning mostly preserves that domain-adapted space.}}\\[1pt]
{\footnotesize Provides a layer-wise mechanistic account of why staged adaptation can work: compares multiple fine-tuning paths and tracks how encoder representations and linearly decodable signals change across layers and checkpoints. Fine-tune Whisper (various sizes) under regimes: English medical, German diagnostic (low data), direct EN+DE medical, and two-stage EN-medical then EN+DE. Measure layer-wise representation drift across checkpoints and train linear probes for language, domain, and error-predictive cues; relate to WER. In the two-stage recipe, English medical tuning causes the dominant encoder shift; multilingual continuation is a smaller adjustment that largely retains the domain-adapted representation. Language/domain remain decodable, while error-predictive linear cues weaken as WER improves, supporting staged training as ``domain surgery'' followed by light multilingual alignment.}

\medskip


\section*{LLM Evaluation}

{\small\textbf{Judge, Retrieve, or Abstain: Uncertainty-Guarded LLM Judging with Provable Risk Guarantees}} {\scriptsize[\href{https://arxiv.org/abs/2608.17994}{arxiv}]}\\
{\scriptsize Sher Badshah, Ali Emami, Hassan Sajjad}\\

{\small\textit{LLM judging can be made risk-controlled: accept, abstain, or retrieve evidence with finite-sample guarantees on the error rate of accepted verdicts.}}\\[1pt]
{\footnotesize Recasts LLM-as-judge as a calibrated acceptance problem controlling error among accepted decisions (FDR-style) at user-chosen level α, and extends it to a two-stage ``judge → retrieve if uncertain'' pipeline while preserving guarantees. On a labeled calibration set, choose uncertainty thresholds using finite-sample Clopper--Pearson bounds so accepted judgments have error ≤ α with high probability. At test time, accept low-uncertainty parametric judgments; otherwise run retrieval-augmented judging and apply a second calibrated threshold or abstain. Empirically meets the promised reliability target across conditions, and two-stage routing increases coverage: many cases that would be abstained by the parametric judge become safely decidable after retrieval, without violating the calibrated bound on accepted mistakes.}

\medskip



\section*{Field Overview}
{\footnotesize Today's submissions are dominated by agentic LLMs and how to evaluate/control them: at least \textasciitilde{}40 papers on LLM agents, multi-agent orchestration, tool use, harness optimization, credit assignment, and ``LLM-as-a-judge'' reliability (e.g., agent benchmarks like Thinkingbox/ComponentBench/DeltaML-Bench, judge routing/abstention, step-level credit audits, skill selection/skill transfer). A second major cluster is long-context efficiency + memory (≈15--20 papers) spanning KV-cache reuse/compression (FlashPrefill V2, ReCache, VoxZip), sparse/alternative attention, modular neural memory (MoNe/Proteus/ArborMem/MemFuse), and new memory failure benchmarks (MemTrapBench, memory commitment/clarification studies). Safety/security/privacy is also heavily represented (≈15+): unlearning (ConceptGuard + GRPO-unlearning reliability), jailbreak/DoS attacks (TempJail, SMTrap, Never Stop Speaking), watermarking robustness/fairness, context leakage, split-learning gradients, and prompt-safety guardrails.

A striking surprise is the sheer volume of audio/speech/music work---well over 150 titles---covering audio-language models, TTS/voice agents, audio deepfake detection/watermarking, long-form/streaming ASR, and text-to-music/audio generation benchmarks (e.g., AT-ADD/MADBench/MMAC/AudioScape-TTA, multiple Qwen-Audio reports, many TTS controllability and low-latency/full-duplex systems). Domain-heavy evaluation is another visible direction (≈20): medicine (HealMed, MedUAG, radiology agents, PHI recovery), legal (ContractScrub, GreekBarRetrieval, InsufficiencyBench, CoAL-RAG), finance (NER under shift, temporal leakage audits, tariff recommendation), and public-sector/local-language benchmarking---suggesting a trend toward ``real deployment conditions'' (multilinguality, distribution shift, executable environments) rather than generic leaderboards.}


\end{document}